% Part: first-order-logic
% Chapter: lindstrom
% Section: lindstrom-proof

\documentclass[../../../include/open-logic-section]{subfiles}

\begin{document}

\olfileid[pt]{mod}{lin}{prf}

\olsection{Teorema de Lindstr\"om}

\begin{lem}
\ollabel{lem:lindstrom}
Suponha que $!E \in L(\Lang{L})$, com $\Lang{L}$ finita, e suponha
também que exista $n \in \Nat$ tal que, para quaisquer duas
!!{structure}s $\Struct{M}$ e~$\Struct{N}$, se
$\Struct{M} \equiv_n \Struct{N}$ e $\Struct{M} \models_L !E$, então
também $\Struct{N} \models_L !E$. Nesse caso, $!E$ é equivalente a uma
!!{sentence} de primeira ordem, O.s.\@, existe um $!D$ de primeira
ordem tal que $\Mod(L){!E} = \Mod(L){!D}$.
\end{lem}

\begin{proof} 
Seja $n$ tal que quaisquer duas !!{structure}s $n$-equivalentes
$\Struct{M}$ e $\Struct{N}$ concordem quanto ao valor atribuído
a~$!E$. Lembre \olref[bas][pis]{prop:qr-finite}: a menos de
equivalência lógica, há apenas um número finito de !!{sentence}s de
primeira ordem, em uma !!{language} finita, com posto de quantificadores
não maior que~$n$. Agora, para cada !!{structure}~$\Struct{M}$ fixa,
seja $!D_{\Struct{M}}$ a conjunção de todas as !!{sentence}s de
primeira ordem~$!E$ verdadeiras em~$\Struct{M}$ com
$\QuantRank{!E} \le n$ (essa conjunção é finita), de modo que
$\Struct{N} \models !D_{\Struct{M}}$ se, e somente se,
$\Struct{N} \equiv_n \Struct{M}$. Tome então
$!D = \textstyle\bigvee
\Setabs{!D_{\Struct{M}}}{\Struct{M} \models_L !E}$; essa disjunção
também é finita (a menos de equivalência lógica).

Segue-se a conclusão $\Mod(L){!E} = \Mod(L){!D}$. De fato, se
$\Struct{N} \models_L !D$, então, para alguma
$\Struct{M} \models_L !E$, temos
$\Struct{N} \models !D_{\Struct{M}}$, donde também
$\Struct{N} \models_L !E$ (pela hipótese do lema). Reciprocamente, se
$\Struct{N} \models_L !E$, então $!D_\Struct{N}$ é um disjunto em
$!D$; como $\Struct{N} \models !D_\Struct{N}$, também
$\Struct{N} \models_L !D$.
\end{proof}

\begin{thm}[Teorema de Lindstr\"om]
  \ollabel{thm:lindstrom} Suponha que $\tuple{L, \models_L}$ tenha as
  Propriedades da Compacidade e de L\"owenheim--Skolem. Então
  $\tuple{L, \models_L} \le \tuple{F, \models}$ (portanto,
  $\tuple{L, \models_L}$ é equivalente à lógica de primeira ordem).
\end{thm}

\begin{proof}
Por \olref{lem:lindstrom}, basta mostrar que, para todo
$!E \in L(\Lang{L})$, com $\Lang{L}$ finita, existe $n \in \Nat$ tal
que, para quaisquer duas !!{structure}s $\Struct{M}$ e~$\Struct{N}$:
se $\Struct{M} \equiv_n \Struct{N}$, então $\Struct{M}$ e $\Struct{N}$
concordam quanto a~$!E$. Com efeito, nesse caso $!E$ é equivalente a
uma !!{sentence} de primeira ordem, donde
$\tuple{L, \models_L} \le \tuple{F, \models}$. Como trabalhamos em uma
!!{language} finita e puramente relacional, por
\olref[bas][pis]{thm:b-n-f} podemos substituir a afirmação
$\Struct{M} \equiv_n \Struct{N}$ pela afirmação algébrica
correspondente $I_n(\emptyset,\emptyset)$.

Dado $!E$, suponha, para obter uma contradição, que para cada $n$
existam !!{structure}s $\Struct{M}_n$ e $\Struct{N}_n$ tais que
$I_n(\emptyset, \emptyset)$, mas (digamos)
$\Struct{M}_n \models_L !E$ enquanto
$\Struct{N}_n \not\models_L !E$. Pela Propriedade do Isomorfismo,
podemos supor que todas as $\Struct{M}_n$ interpretam as constantes da
linguagem pelos mesmos objetos; além disso, como há apenas um número
finito de !!{sentence}s atômicas na linguagem, podemos supor também
que elas satisfazem as mesmas !!{sentence}s atômicas (caso contrário,
podemos tomar uma subsequência das $\Struct{M}_n$). Seja $\Struct{M}$ a
união de todas as $\Struct{M}_n$, O.s.\@, a única !!{structure} mínima
que tem cada $\Struct{M}_n$ como subestrutura. Como na demonstração de
\olref[lsp]{thm:abstract-p-isom}, seja $\Struct{M}^*$ a extensão de
$\Struct{M}$ com !!{domain}
$\Domain{M} \cup \Domain{M}^{<\omega}$, na !!{language} expandida que
compreende os predicados de concatenação $P$ e~$Q$.

De modo semelhante, defina $\Struct{N}_n$, $\Struct{N}$ e
$\Struct{N}^*$. Agora, seja $\Struct{M}$ a !!{structure} cujo
!!{domain} compreende os !!{domain}s de $\Struct{M}^*$ e
$\Struct{N}^*$, bem como os números naturais~$\Nat$ com sua ordenação
natural~$\le$, na !!{language} com predicados adicionais que
representam os !!{domain}s $\Domain{M}$, $\Domain{N}$,
$\Domain{M}^{<\omega}$ e $\Domain{N}^{<\omega}$, além de predicados que
codificam os domínios de $\Struct{M}_n$ e $\Struct{N}_n$ no sentido de
que:
\begin{align*}
  \Domain{M_n} & = \Setabs{a \in \Domain{M}}{R(a, n)}; & 
  \Domain{N_n} & = \Setabs{a \in \Domain{N}}{S(a,n)}; \\
  \Domain{M}^{<\omega}_n & = \Setabs{a \in \Domain{M}^{<\omega}}{R(a,n)}; &
  \Domain{N}^{<\omega}_n & = \Setabs{a \in \Domain{N}^{<\omega}}{S(a,n)}. 
\end{align*}
A !!{structure}~$\Struct{M}$ também tem uma relação ternária $J$ tal
que $J(n, \mathbf{a}, \mathbf{b})$ vale se, e somente se,
$I_n(\mathbf{a}, \mathbf{b})$.

Existe agora !!a{sentence}~$!D$ na !!{language}~$\Lang{L}$ ampliada
com $R$, $S$, $J$ etc., que afirma que $\le$ é uma ordenação linear
discreta com primeiro elemento, mas sem último elemento, e que
$\Struct{M}_n \models !E$, $\Struct{N}_n \not\models !E$ e, para cada
$n$ na ordenação, $J(n, \mathbf{a}, \mathbf{b})$ vale se, e somente
se, $I_n(\mathbf{a}, \mathbf{b})$.

Usando a Propriedade da Compacidade, podemos encontrar um modelo
$\Struct{M}^*$ de $!D$ no qual a ordenação contém um elemento não
padrão~$n^*$. Em particular, $\Struct{M^*}$ conterá então
sub!!{structure}s $\Struct{M_{n^*}}$ e $\Struct{N_{n^*}}$ tais que
$\Struct{M_{n^*}} \models_L !E$ e
$\Struct{N_{n^*}} \not\models_L !E$. Mas agora podemos definir um
conjunto $\mathcal{I}$ de pares de $k$-uplas de $\Domain{M_{n^*}}$ e
$\Domain{N_{n^*}}$, tomando
$\tuple{\mathbf{a}, \mathbf{b}} \in \mathcal{I}$ se, e somente se,
$J(n^*-k, \mathbf{a}, \mathbf{b})$, onde $k$ é o comprimento de
$\mathbf{a}$ e $\mathbf{b}$. Como $n^*$ é não padrão, para todo $k$
padrão temos $n^* - k >0$, e o conjunto $\mathcal{I}$ testemunha que
$\Struct{M_{n^*}} \simeq_p \Struct{N_{n^*}}$. Mas, por
\olref[lsp]{thm:abstract-p-isom}, $\Struct{M_{n^*}}$ é $L$-equivalente
a $\Struct{N_{n^*}}$, uma contradição.
\end{proof}

\end{document}
